\section{Introduction}
\label{sec:introduction}

Long-lived sensor deployments operate under a basic mismatch: the environment
is bursty, but their schedules are static.  A bridge can remain quiescent for
hours before a short vibration transient; indoor air quality changes slowly
until a door opens; and soil moisture drifts until rain produces a step.  A
deployment configured for the transient repeatedly measures predictable steady
state.  One configured for the steady state discovers the event too late.
This tension is especially costly on energy-harvesting and hard-to-service
nodes, where a sampling decision consumes not only sensor energy but also CPU
wake-up and radio-tail energy~\cite{lorincz2004sensor,lin2010energy}.

Adaptive sampling is not new.  Prior systems use prediction error, compressive
sensing, event triggers, or application-specific rules to vary a sampling
rate~\cite{deshpande2004model,raghunathan2006energyscavenging,jain2004duality,lu2010smart}.  Three practical gaps remain.  First, a residual is meaningful only when paired with the model's uncertainty: a
2$^\circ$ error can be routine during a volatile afternoon and alarming at
night.  Second, a locally sensible rate can still exhaust a finite daily energy
budget after an extended event.  Third, correlated nodes often make the same
decision and report redundant evidence.  Existing techniques commonly address
one gap while treating the others as configuration problems.

We ask a narrower systems question: \emph{can a node allocate a fixed energy
budget to the moments at which one additional observation is most likely to
matter?}  \system answers with a closed-loop scheduler whose state fits in 184
bytes and whose update requires 1.7k integer-equivalent operations.  It
combines a constant-velocity predictor with an exponentially weighted error
model.  The resulting risk score controls a continuous sampling interval.  A
token bucket stretches that interval when the node spends energy too quickly,
while hard age and change-point guards retain responsiveness.  Finally, a
four-byte neighbor digest discounts observations that are likely to be
spatially redundant.

The design intentionally avoids a heavyweight learned policy.  This makes the
decision path inspectable, permits safe bounds on sample age, and lets one set
an energy budget without retraining.  The trade-off is that \system is a
scheduler, not a semantic event classifier: it decides \emph{when} to observe,
then invokes the deployment's existing detector.

This paper makes three contributions:

\begin{itemize}[leftmargin=*,itemsep=2pt,topsep=3pt]
  \item We formulate sampling as online, risk-weighted budget allocation and
  derive a bounded interval controller that combines uncertainty, surprise,
  staleness, and spatial redundancy (\S\ref{sec:design}).
  \item We design a budget governor and neighbor protocol that tolerate energy
  droughts, event storms, loss, and cold starts without requiring a gateway in
  the control loop (\S\ref{sec:implementation}).
  \item We provide a trace-driven evaluation spanning five signal regimes and
  report energy, event recall, delay, traffic, sensitivity, and failure modes
  (\S\ref{sec:evaluation}).
\end{itemize}

All results in this manuscript are from deterministic trace replay and the
hardware energy model in Table~\ref{tab:platform}; they are not presented as a
field deployment.  This distinction makes the evidence reproducible and keeps
the conclusions scoped to scheduling behavior.
